\section{Граф знаний и интерпретация рекомендаций}

\subsection{Роль слоя интерпретации в работе}

Ранжированный список вероятностей сам по себе ещё не даёт содержательного объяснения,
почему именно эти кандидаты оказались в верхней части списка и какие факторы
поддерживают или ослабляют отдельные варианты. В рекомендательной постановке
пользователь получает не один ответ, а короткий список вероятных следующих культур,
поэтому поверх модельного результата вводится отдельный слой интерпретации на основе
графа знаний.

В данной работе используется компактный граф знаний предметной области. Его задача состоит в интерпретации
кандидатов, уже сформированных предиктивной моделью, через переходную статистику,
региональную типичность, принадлежность к группам культур и срабатывание поддерживающих
или предупреждающих правил.

Такой подход позволяет избежать смешения двух разных функций. Вероятностный вывод модели
и его семантическая интерпретация в работе разведены по разным уровням. Графовый слой
показывает, насколько тот или иной кандидат согласован с наблюдаемым контекстом.

\subsection{Общая структура графа знаний}

Граф знаний в работе построен как компактное графовое представление сущностей,
связанных с задачей рекомендации следующей культуры. По метаданным запуска граф
содержит 2827 узлов и 18154 ребра, а также включает 5 интерпретационных правил и
набор разборов примеров, используемых для анализа типичных сценариев поведения системы.

В графе используются четыре типа сущностей:

\begin{itemize}
\tightlist
\item
  культуры;
\item
  группы культур;
\item
  регионы;
\item
  правила.
\end{itemize}

Такая структура позволяет представить не только сами классы целевых культур, но и их
более общий контекст. Культуры связываются между собой переходными отношениями,
дополнительно привязываются к укрупнённым группам, рассматриваются в региональном
контексте и получают интерпретационные сигналы через набор правил. Тем самым граф
знаний выступает как слой предметно-ориентированного представления, который можно
использовать для анализа списка кандидатов.

В самом общем виде место графа в архитектуре можно описать как последовательность:
\begin{center}\(
\text{список кандидатов модели}
\rightarrow
\text{контекст графа и правил}
\rightarrow
\text{пояснение для кандидата}.
\)\end{center}

Эта логика подчёркивает, что граф не порождает кандидатов сам по себе, а работает
с уже полученным списком и дополняет его интерпретационным контекстом.

\subsection{Типы связей в графе}

Структура графа опирается на несколько видов связей. В первую очередь это переходы
между культурами, построенные по статистике, рассчитанной только на обучающей выборке.
Они отражают наблюдаемую частоту перехода от одной культуры к другой и позволяют
использовать накопленные последовательности как источник поддерживающего сигнала для
кандидатов. Кроме того, в граф включены связи принадлежности культуры к укрупнённой
группе, а также связи, отражающие типичность культуры для региона. Наконец, отдельный
тип связей используется для связывания интерпретационных правил с конкретными
кандидатами.

Смысл каждого из этих типов связей различен. Переходные связи показывают, насколько
кандидат согласуется с наблюдаемыми историческими закономерностями. Связи с группами
культур позволяют анализировать более общий уровень сходства и насыщения. Региональные
связи дают дополнительный сигнал о том, является ли культура типичной для
рассматриваемого контекста. Связи с правилами фиксируют поддерживающие и
предупреждающие эффекты, которые используются в логике интерпретации. В совокупности
это позволяет не сводить анализ к одной числовой вероятности, а рассматривать
кандидата в более богатом контексте.

\subsection{Интерпретационные правила}

Отдельный компонент слоя интерпретации образует набор правил. В работе используются
пять правил; их смысл и тип интерпретационного сигнала приведены в
таблице~\ref{tab:kg-rules}.

\begin{table}[H]
\centering
\small
\caption{Интерпретационные правила knowledge graph}
\label{tab:kg-rules}
\begin{tabular}{>{\raggedright\arraybackslash}p{0.28\textwidth}
                >{\raggedright\arraybackslash}p{0.48\textwidth}
                >{\raggedright\arraybackslash}p{0.16\textwidth}}
\toprule
Правило & Смысл & Тип сигнала \\
\midrule
\texttt{repeat\_crop\_warning}
  & предупреждение, если кандидат повторяет последнюю культуру в истории поля
  & warning \\
\texttt{same\_group\_saturation}
  & предупреждение, если история насыщена одной группой культур и кандидат относится к той же группе
  & warning \\
\texttt{legume\_break\_bonus}
  & поддерживающий сигнал для кандидатов из бобовых культур, если в недавней истории не было бобовых
  & support \\
\texttt{fallow\_break\_effect}
  & поддержка возврата из \texttt{fallow} к культуре или предупреждение при повторном выборе \texttt{fallow}
  & support / warning \\
\texttt{region\_typicality}
  & поддерживающий сигнал для культур, типичных для региона по train-only статистике
  & support \\
\bottomrule
\end{tabular}

\vspace{0.4em}
\footnotesize Источник: \path{artifacts/results/knowledge_graph_explainability/kg_rules.csv}.
\normalsize
\end{table}


Содержательно правила можно разделить на два типа. Первый тип связан с
предупреждениями: повтор одной и той же культуры или насыщение последовательности
одной группой культур могут рассматриваться как предупреждающие сигналы. Второй тип
связан с поддерживающими сигналами: например, логика разрыва севооборота за счёт
бобовых, эффект пара или региональная типичность культуры могут усиливать
интерпретационную поддержку отдельного кандидата. Такое сочетание позволяет
анализировать список кандидатов не только с точки зрения статистического ранжирования,
но и с точки зрения более понятной семантической логики.

\subsection{Интерпретация на уровне отдельных кандидатов}

Слой интерпретации в работе применяется к отдельным позициям ранжированного списка.
Граф не предназначен для объяснения «внутренностей» CatBoost
напрямую и не интерпретирует веса модели в терминах причинности. Вместо этого для
каждой позиции в списке из \(k\) кандидатов собирается набор дополнительных
характеристик, описывающих его положение в контексте графа и правил.

Для каждого кандидата формируются следующие элементы интерпретационного вывода:

\begin{itemize}
\tightlist
\item
  положение кандидата в ранжированном списке;
\item
  вероятность или оценка модели;
\item
  принадлежность к укрупнённой группе;
\item
  переходная поддержка;
\item
  региональный сигнал;
\item
  поддерживающие и предупреждающие флаги;
\item
  итоговый \texttt{graph\_support\_score};
\item
  уровень графовой поддержки.
\end{itemize}

Тем самым выход слоя интерпретации можно описать так:

\begin{center}\(
\text{кандидат}
\rightarrow
\text{оценка модели}
+
\text{контекст графа и правил}
+
\text{краткое пояснение}
\)\end{center}

Такой формат удобен по двум причинам. Во-первых, он сохраняет связь с модельным
результатом: каждый кандидат остаётся частью исходного ранжированного списка.
Во-вторых, он даёт более содержательный комментарий к каждой позиции в списке.
Пользователь видит не только сам факт того, что культура попала в список из
\(k\) кандидатов, но и дополнительный контекст, объясняющий, какие сигналы её
поддерживают или ослабляют.

\subsection{Показатель графовой поддержки и его интерпретация}

\texttt{graph\_support\_score} агрегирует переходную статистику,
региональную типичность, поддержку правил и предупреждения в единую
оценку того, насколько кандидат согласован с контекстом графа.
По своей природе это интерпретационный показатель: он не является
вероятностью и не претендует на роль альтернативного ранжирования —
модельные вероятности CatBoost остаются основой списка кандидатов.
Полезность \texttt{graph\_support\_score} в другом: он позволяет
сопоставлять кандидатов по степени их семантической поддержки и
выделять случаи, где первая рекомендация модели расходится
с интерпретационным контекстом.

\subsection{Согласованность между моделью и интерпретационным слоем}

Для анализа поведения списка кандидатов в работе вводится понятие согласованности
между первым выбором модели и контекстом графа и правил. Это позволяет рассматривать
не только положение кандидатов в ранжировании, но и то, насколько первый кандидат
модели поддержан слоем интерпретации. По итогам разборов примеров были получены
случаи с \texttt{strong} и \texttt{weak\_alignment}; категория \texttt{medium}
в данном запуске не появлялась, поскольку выбранные примеры были отнесены только
к крайним типам согласованности.

Смысл этой интерпретации заключается не в делении на «правильные» и «неправильные»
ответы графа. Например, \texttt{weak\_alignment} не означает, что модель ошиблась,
а граф оказался прав. Он означает более узкую вещь: выбранный первый кандидат модели
получает относительно слабую поддержку со стороны контекста графа и правил, тогда как
некоторые альтернативы из списка кандидатов могут выглядеть семантически более
поддержанными. В прикладной постановке это особенно важно, поскольку показывает,
в каких случаях список кандидатов следует рассматривать внимательнее, а не
ограничиваться первой позицией.

За счёт этого графовый слой помогает отделять согласованные случаи
от ситуаций с конкурирующими альтернативами.
Это усиливает всю рекомендательную постановку работы, поскольку показывает, что
полезность списка кандидатов заключается не только в наличии нескольких кандидатов,
но и в возможности содержательно сравнивать их между собой.

\subsection{Место графа знаний в общей логике работы}

В общей архитектуре системы граф знаний располагается после рекомендательного слоя
и слоя оценки уверенности. Сначала модель формирует список наиболее вероятных культур.
Затем слой оценки уверенности показывает, насколько надёжно этот список выглядит
с точки зрения вероятностного вывода. После этого граф знаний даёт пояснение на
уровне отдельного кандидата. Такая последовательность важна концептуально: сначала
формируется статистический результат, затем оценивается его надёжность, и только
после этого добавляется интерпретационный контекст.

Благодаря такому расположению граф знаний дополняет основную модель и делает результат
менее закрытым для анализа. Пользователь получает не только ранжированный список культур,
но и пояснения, связанные с переходными закономерностями, региональным контекстом и интерпретационными правилами.
Такое разделение ролей делает интерпретацию рекомендаций более понятной.

Таким образом, граф знаний дополняет вероятностное ранжирование основной модели
и используется для интерпретации уже сформированного списка кандидатов. Его
назначение состоит в добавлении предметного контекста к модельному результату:
переходной поддержки, региональной типичности, принадлежности к группам культур
и предупреждающих сигналов. Особенно полезным такой слой оказывается в спорных
случаях, когда первая рекомендация модели не является единственным правдоподобным
вариантом. Далее работа графового слоя рассматривается на отдельных примерах,
где можно сопоставить историю поля, вероятности модели, уровень уверенности и
интерпретационный контекст для конкретных кандидатов.
